\documentclass{amsart}
% For dealing with references we use the comment environment
\usepackage{verbatim}
\newenvironment{reference}{\comment}{\endcomment}
%\newenvironment{reference}{}{}
\newenvironment{slogan}{\comment}{\endcomment}
\newenvironment{history}{\comment}{\endcomment}

% For commutative diagrams we use Xy-pic
\usepackage[all]{xy}

% We use 2cell for 2-commutative diagrams.
\xyoption{2cell}
\UseAllTwocells

% We use multicol for the list of chapters between chapters
\usepackage{multicol}

% This is generally recommended for better output
\usepackage{lmodern}
\usepackage[T1]{fontenc}

% For cross-file-references

% Package for hypertext links:
\usepackage{hyperref}

% For any local file, say "hello.tex" you want to link to please

% Theorem environments.
%
\theoremstyle{plain}
\newtheorem{theorem}[subsection]{Theorem}
\newtheorem{proposition}[subsection]{Proposition}
\newtheorem{lemma}[subsection]{Lemma}

\theoremstyle{definition}
\newtheorem{definition}[subsection]{Definition}
\newtheorem{example}[subsection]{Example}
\newtheorem{exercise}[subsection]{Exercise}
\newtheorem{situation}[subsection]{Situation}

\theoremstyle{remark}
\newtheorem{remark}[subsection]{Remark}
\newtheorem{remarks}[subsection]{Remarks}

\numberwithin{equation}{subsection}

% Macros
%
\def\lim{\mathop{\mathrm{lim}}\nolimits}
\def\colim{\mathop{\mathrm{colim}}\nolimits}
\def\Spec{\mathop{\mathrm{Spec}}}
\def\Hom{\mathop{\mathrm{Hom}}\nolimits}
\def\Ext{\mathop{\mathrm{Ext}}\nolimits}
\def\SheafHom{\mathop{\mathcal{H}\!\mathit{om}}\nolimits}
\def\SheafExt{\mathop{\mathcal{E}\!\mathit{xt}}\nolimits}
\def\Sch{\mathit{Sch}}
\def\Mor{\mathop{\mathrm{Mor}}\nolimits}
\def\Ob{\mathop{\mathrm{Ob}}\nolimits}
\def\Sh{\mathop{\mathit{Sh}}\nolimits}
\def\NL{\mathop{N\!L}\nolimits}
\def\CH{\mathop{\mathrm{CH}}\nolimits}
\def\proetale{{pro\text{-}\acute{e}tale}}
\def\etale{{\acute{e}tale}}
\def\QCoh{\mathit{QCoh}}
\def\Ker{\mathop{\mathrm{Ker}}}
\def\Im{\mathop{\mathrm{Im}}}
\def\Coker{\mathop{\mathrm{Coker}}}
\def\Coim{\mathop{\mathrm{Coim}}}

% Boxtimes
%
\DeclareMathSymbol{\boxtimes}{\mathbin}{AMSa}{"02}

%
% Macros for moduli stacks/spaces
%
\def\QCohstack{\mathcal{QC}\!\mathit{oh}}
\def\Cohstack{\mathcal{C}\!\mathit{oh}}
\def\Spacesstack{\mathcal{S}\!\mathit{paces}}
\def\Quotfunctor{\mathrm{Quot}}
\def\Hilbfunctor{\mathrm{Hilb}}
\def\Curvesstack{\mathcal{C}\!\mathit{urves}}
\def\Polarizedstack{\mathcal{P}\!\mathit{olarized}}
\def\Complexesstack{\mathcal{C}\!\mathit{omplexes}}
% \Pic is the operator that assigns to X its picard group, usage \Pic(X)
% \Picardstack_{X/B} denotes the Picard stack of X over B
% \Picardfunctor_{X/B} denotes the Picard functor of X over B
\def\Pic{\mathop{\mathrm{Pic}}\nolimits}
\def\Picardstack{\mathcal{P}\!\mathit{ic}}
\def\Picardfunctor{\mathrm{Pic}}
\def\Deformationcategory{\mathcal{D}\!\mathit{ef}}

\setlength{\emergencystretch}{3em}
\begin{document}
\title[Stacks Project, Tag 02S9]{447 --- Dimension-filtered coherent Grothendieck groups identify with cycle groups}
\maketitle
\noindent Copyright (C) 2005--2025 Johan de Jong. Stacks Project authors. Source: \href{https://stacks.math.columbia.edu/tag/02S9}{Tag 02S9}.

\noindent Editorial difficulty note (not author text). Difficulty H2: The proof must recover the coherent class from all generic-point lengths, including infinitely many locally finite supports. It globalizes generic composition series, constructs coherent direct sums of successive filtration layers, and compares these with structure sheaves by extended ideal subsheaves whose cokernels have lower-dimensional support..

\noindent Editorial provenance note (not author text). History: excerpt from revision \texttt{a0444\allowbreak 6e57e\allowbreak c1fbc\allowbreak 252a8\allowbreak 71afc\allowbreak ec775\allowbreak 2fb28\allowbreak 07b14}, chow.tex, lines 3756--3895. Reference presentation was adapted; original-author context and core support blocks were retained or added. The mathematical wording is unchanged. Licensed under GNU FDL 1.2 or later; no invariant sections or cover texts. See ../licenses/COPYING. Context blocks reproduce author definitions and supporting results; they are not additional counted problems.

\section{Chow groups and K-groups}
\label{section-chow-and-K}

\noindent
In this section we are going to compare $K_0$ of the
category of coherent sheaves to the chow groups.

\medskip\noindent
Let $(S, \delta)$ be as in Situation \ref{situation-setup}.
Let $X$ be a scheme locally of finite type over $S$.
We denote $\textit{Coh}(X) = \textit{Coh}(\mathcal{O}_X)$
the category of coherent sheaves on $X$.
It is an abelian category, see
Cohomology of Schemes, Lemma \href{https://stacks.math.columbia.edu/tag/01Y0}{lemma coherent abelian Noetherian (Tag 01Y0)}.
For any $k \in \mathbf{Z}$ we let $\textit{Coh}_{\leq k}(X)$
be the full subcategory of $\textit{Coh}(X)$
consisting of those coherent sheaves $\mathcal{F}$
having $\dim_\delta(\text{Supp}(\mathcal{F})) \leq k$.



\begin{situation}
\label{situation-setup}
Here $S$ is a locally Noetherian, and universally catenary scheme.
Moreover, we assume $S$ is endowed with a dimension function
$\delta : S \longrightarrow \mathbf{Z}$.
\end{situation}

\begin{lemma}
\label{lemma-cycles-k-group}
Let $(S, \delta)$ be as in Situation \ref{situation-setup}.
Let $X$ be a scheme locally of finite type over $S$.
The maps
$$
Z_k(X)
\longrightarrow
K_0(\textit{Coh}_{\leq k}(X)/\textit{Coh}_{\leq k - 1}(X)),
\quad
\sum n_Z[Z] \mapsto
\left[\bigoplus\nolimits_{n_Z > 0} \mathcal{O}_Z^{\oplus n_Z}\right]
-
\left[\bigoplus\nolimits_{n_Z < 0} \mathcal{O}_Z^{\oplus -n_Z}\right]
$$
and
$$
K_0(\textit{Coh}_{\leq k}(X)/\textit{Coh}_{\leq k - 1}(X))
\longrightarrow
Z_k(X),\quad
\mathcal{F} \longmapsto [\mathcal{F}]_k
$$
are mutually inverse isomorphisms.
\end{lemma}

\begin{proof}
Note that if $\sum n_Z[Z]$ is in $Z_k(X)$, then
the direct sums
$\bigoplus\nolimits_{n_Z > 0} \mathcal{O}_Z^{\oplus n_Z}$ and
$\bigoplus\nolimits_{n_Z < 0} \mathcal{O}_Z^{\oplus -n_Z}$
are coherent sheaves on $X$ since the family $\{Z \mid n_Z > 0\}$
is locally finite on $X$.
The map $\mathcal{F} \to [\mathcal{F}]_k$ is additive
on $\textit{Coh}_{\leq k}(X)$, see
Lemma \href{https://stacks.math.columbia.edu/tag/02QZ}{lemma additivity sheaf cycle (Tag 02QZ)}. And $[\mathcal{F}]_k = 0$
if $\mathcal{F} \in \textit{Coh}_{\leq k - 1}(X)$. By part (1)
of Homology, Lemma \href{https://stacks.math.columbia.edu/tag/02MX}{lemma serre subcategory K groups (Tag 02MX)}
this implies that the second map is well defined too.
It is clear that the composition of the first map with the second
map is the identity.

\medskip\noindent
Conversely, say we start with a coherent sheaf $\mathcal{F}$
on $X$. Write $[\mathcal{F}]_k = \sum_{i \in I} n_i[Z_i]$
with $n_i > 0$ and $Z_i \subset X$, $i \in I$
pairwise distinct integral closed subschemes of $\delta$-dimension $k$.
We have to show that
$$
[\mathcal{F}] = [\bigoplus\nolimits_{i \in I} \mathcal{O}_{Z_i}^{\oplus n_i}]
$$
in $K_0(\textit{Coh}_{\leq k}(X)/\textit{Coh}_{\leq k - 1}(X))$.
Denote $\xi_i \in Z_i$ the generic point.
If we set
$$
\mathcal{F}' = \Ker(\mathcal{F} \to \bigoplus \xi_{i, *}\mathcal{F}_{\xi_i})
$$
then $\mathcal{F}'$ is the maximal coherent submodule of $\mathcal{F}$
whose support has dimension $\leq k - 1$. In particular $\mathcal{F}$
and $\mathcal{F}/\mathcal{F}'$ have the same class in
$K_0(\textit{Coh}_{\leq k}(X)/\textit{Coh}_{\leq k - 1}(X))$.
Thus after replacing $\mathcal{F}$ by $\mathcal{F}/\mathcal{F}'$
we may and do assume that the kernel $\mathcal{F}'$ displayed
above is zero.

\medskip\noindent
For each $i \in I$ we choose a filtration
$$
\mathcal{F}_{\xi_i} = \mathcal{F}_i^0 \supset \mathcal{F}_i^1 \supset
\ldots \supset \mathcal{F}_i^{n_i} = 0
$$
such that the successive quotients are of dimension $1$ over the residue
field at $\xi_i$. This is possible as the length of $\mathcal{F}_{\xi_i}$
over $\mathcal{O}_{X, \xi_i}$ is $n_i$.
For $p > n_i$ set $\mathcal{F}_i^p = 0$. For $p \geq 0$ we denote
$$
\mathcal{F}^p =
\Ker\left(\mathcal{F} \longrightarrow \bigoplus
\xi_{i, *}(\mathcal{F}_{\xi_i}/\mathcal{F}_i^p)\right)
$$
Then $\mathcal{F}^p$ is coherent, $\mathcal{F}^0 = \mathcal{F}$, and
$\mathcal{F}^p/\mathcal{F}^{p + 1}$ is isomorphic to a free
$\mathcal{O}_{Z_i}$-module of rank $1$ (if $n_i > p$) or $0$
(if $n_i \leq p$) in an open neighbourhood of $\xi_i$. Moreover,
$\mathcal{F}' = \bigcap \mathcal{F}^p = 0$. Since every quasi-compact
open $U \subset X$ contains only a finite number of $\xi_i$
we conclude that $\mathcal{F}^p|_U$ is zero for $p \gg 0$.
Hence $\bigoplus_{p \geq 0} \mathcal{F}^p$ is a coherent
$\mathcal{O}_X$-module. Consider the short exact sequences
$$
0 \to
\bigoplus\nolimits_{p > 0} \mathcal{F}^p \to
\bigoplus\nolimits_{p \geq 0} \mathcal{F}^p \to
\bigoplus\nolimits_{p > 0} \mathcal{F}^p/\mathcal{F}^{p + 1} \to 0
$$
and
$$
0 \to
\bigoplus\nolimits_{p > 0} \mathcal{F}^p \to
\bigoplus\nolimits_{p \geq 0} \mathcal{F}^p \to
\mathcal{F} \to 0
$$
of coherent $\mathcal{O}_X$-modules. This already shows that
$$
[\mathcal{F}] = [\bigoplus \mathcal{F}^p/\mathcal{F}^{p + 1}]
$$
in $K_0(\textit{Coh}_{\leq k}(X)/\textit{Coh}_{\leq k - 1}(X))$.
Next, for every $p \geq 0$ and $i \in I$ such that $n_i > p$
we choose a nonzero ideal sheaf $\mathcal{I}_{i, p} \subset \mathcal{O}_{Z_i}$
and a map $\mathcal{I}_{i, p} \to \mathcal{F}^p/\mathcal{F}^{p + 1}$ on $X$
which is an isomorphism over the open neighbourhood of $\xi_i$
mentioned above. This is possible by
Cohomology of Schemes, Lemma \href{https://stacks.math.columbia.edu/tag/0FD0}{lemma extend coherent (Tag 0FD0)}.
Then we consider the short exact sequence
$$
0 \to
\bigoplus\nolimits_{p \geq 0, i \in I, n_i > p} \mathcal{I}_{i, p}
\to
\bigoplus \mathcal{F}^p/\mathcal{F}^{p + 1} \to
\mathcal{Q} \to 0
$$
and the short exact sequence
$$
0 \to
\bigoplus\nolimits_{p \geq 0, i \in I, n_i > p} \mathcal{I}_{i, p}
\to
\bigoplus\nolimits_{p \geq 0, i \in I, n_i > p} \mathcal{O}_{Z_i}
\to
\mathcal{Q}' \to 0
$$
Observe that both $\mathcal{Q}$ and $\mathcal{Q}'$ are zero in a neighbourhood
of the points $\xi_i$ and that they are supported on $\bigcup Z_i$.
Hence $\mathcal{Q}$ and $\mathcal{Q}'$ are in
$\textit{Coh}_{\leq k - 1}(X)$.
Since
$$
\bigoplus\nolimits_{i \in I} \mathcal{O}_{Z_i}^{\oplus n_i} \cong
\bigoplus\nolimits_{p \geq 0, i \in I, n_i > p} \mathcal{O}_{Z_i}
$$
this concludes the proof.
\end{proof}
\end{document}
